\documentclass[11pt]{article}
\usepackage[a4paper,margin=24mm]{geometry}
\usepackage{amsmath,amssymb,amsthm,booktabs,hyperref}
\newtheorem{theorem}{Theorem}\newtheorem{corollary}{Corollary}
\newcommand{\gen}[1]{\langle#1\rangle}
\title{Transitive string C-groups: the complete half-degree bound}
\author{An attributed application of the common string-separator theory}
\date{12 October 2026}
\begin{document}\maketitle
This application recovers the principal contract of Cameron, Fernandes,
Leemans and Mixer, \emph{String C-groups as transitive subgroups of $S_n$},
arXiv:1410.5863v1, 21 October 2014. Its complete twelve-page edition is
fixed. The complete current Topic B common paper and the supplement
\emph{String separators, normal modules and transitive rank} are the
two declared common inputs. No previous application is required.
All source result numbers and qualifications are recorded in the clause map.
The finite input is a complete exact census, not a sample of markings.

\begin{theorem}[Theorem 1.2, with ordered marking distinctions]
Let $G$ be a finite string C-group of rank $d$ with a faithful transitive
action of degree $n$, with $G\ne A_n,S_n$ in this action.
If $d>n/2$, then one of the following occurs:
\begin{enumerate}
\item $n=2m$ with $m$ odd, and $G=C_2\wr S_m$ has rank $m+1$.
For $m\geq5$, its unordered distinguished set is unique up to marked-group
isomorphism and is $\{\lambda,\lambda+e_1,s_1,\ldots,s_{m-1}\}$,
where $\lambda$ flips all blocks, $e_1$ flips the first block, and
$s_i$ interchanges blocks $i,i+1$ without internal flips.
The two ordered types up to global duality are
$[2,4,3,\ldots,3]$ and $[2,3,\ldots,3,4]$.
\item $G$ is one of the imprimitive degree-six or degree-eight
markings in Table~1 below, with the displayed types and their global reversals.
\item $n=6$, $G\cong S_5$ in its primitive action, and the marking is
the rank-four simplex of type $[3,3,3]$.
\end{enumerate}
The first family also exists at $m=3$ and is included among the
degree-six ordered markings. The residual alternative is $d\leq n/2$.
\end{theorem}
\begin{proof}
Choose a maximal nontrivial block system if the action is imprimitive.
The common block argument gives $d\leq n/2$ when both block size and
number exceed two. For two large blocks of size $k\geq5$, the common
uniform proof and the accepted complete ten-action finite-neighbor lemma
give $d\leq k$. For blocks of size two, the common normal-module argument
gives $d\leq m+1$, with equality precisely as in the theorem. It proves
the actual marking, cocycle normalization and IP, not merely its order.
The complete degree-six/eight/nine census treats all smaller imprimitive
actions; odd prime degrees have no imprimitive action, and the proper
subgroups of $S_4$ have independent length at most two. There is no
rank-five exception in degree nine.

For primitive actions the common sharp bound and exact residual census
give $d<n/2$ except for degree-six $A_5,S_5$ and degree-ten $S_6$ with
the types printed in the common theorem. Only the degree-six rank-four
$S_5$ marking has $d>n/2$. Its Coxeter relations and order $120$ identify
it with the simplex. This proves all branches.
\end{proof}

\begin{table}[ht]\centering
\begin{tabular}{clrl}
\toprule
Degree & Actual group & Order & Ordered types, up to global duality\\
\midrule
6&$S_3\times S_3$&36&$[3,2,3]$\\
6&$C_2\times S_4$&48&$[2,3,3]$\\
6&$C_2\times S_4=C_2\wr S_3$&48&$[2,3,4],\ [2,4,3]$\\
8&$(2^4:S_3):S_3$&576&$[3,4,4,3]$\\
\bottomrule
\end{tabular}
\caption{Complete small imprimitive exception types. Literal permutation
markings and every ordered conjugacy representative are in the required census.}
\end{table}

The source's Table~1 prints $[2,3,3]$ for the order-$36$ group.
That type is impossible: its universal Coxeter group is $C_2\times S_4$
of order $48$, so no quotient has order $36$. The complete exact marking
has type $[3,2,3]$. This is a correction to this fixed source edition;
no claim is made about an unseen later edition.
The source's two written types for the wreath family instead reverse
only its noncentral connected component. They are compatible orders
of the same unordered marked set, but need not be identified as ordered
polytopes or as global duals. The present statement includes both.

The common proof establishes Proposition~2.1 for every $n\geq10$ and
Proposition~2.2 in the standing string-C-group setting of source Section~2:
when $n\geq11$, equality $d=n/2+1$ forces $n\equiv2\pmod4$ and the
unordered pair $(G,S)$ is unique. The complete actual parameter family
shows that this bound is sharp. It does not assert that every transitive
action of the same abstract group has this degree or rank.

\section{The alternating-parabolic consequences}
Write $R$ for the parent rank and $r=R-1$ for a maximal parabolic rank.
The source uses $S=\{\rho_0,\ldots,\rho_d\}$ in Section~4, so there
its printed $d$ is $r$, not $R$.

\begin{corollary}[Source Corollaries 4.1--4.3]
Let $G=A_n$ with a string C-marking of rank $R$, and let $G_i$ be a
maximal parabolic in its natural degree-$n$ action.
If $G_i$ is transitive imprimitive, then $r\leq n/2$, hence
$R\leq\lfloor n/2\rfloor+1$; in particular the source's literal
$r\leq n/2+1$ holds. If $G_i$ is primitive, then $r<n/2$, hence
$R\leq\lceil n/2\rceil$ and the source's literal $r\leq n/2$ holds.
Consequently if $R\geq n/2+2$, every maximal parabolic is intransitive.
\end{corollary}
\begin{proof}
A parabolic is proper by IP, so it is not $A_n$ or $S_n$.
The imprimitive equality family has the odd permutation $\lambda$
when $m$ is odd, and therefore cannot lie in $A_n$. All degree-six
rank-four markings likewise contain an odd permutation; the census
checks this on every actual tuple. The only remaining exception is
the degree-eight rank-five marking in $A_8$. Its parent would have
rank six in $A_8$. The complete rank-six connected-string census in
$A_8$ is empty. Since $A_8$ is directly indecomposable and centerless,
its generating C-diagram must be connected. This rules out that case
and proves $r\leq n/2$.

For a primitive parabolic the degree-six $S_5$ and degree-ten $S_6$
actions are not contained in the corresponding alternating group,
as the explicit permutation models verify. The only remaining
half-degree case is $A_5$ in degree six, of rank three. Its parent
would be a rank-four C-group marking of $A_6$, which the same complete
census excludes. Hence $r<n/2$. The final assertion follows because
$R\geq n/2+2$ would give $r>n/2$.
\end{proof}

The rank convention is essential: in degree five the icosahedral
$A_5$ marking has parent rank three and a transitive primitive
$D_{10}$ parabolic. It satisfies the source's indexed inequality
$r=2\leq5/2$, but not the incorrectly relabelled assertion $R\leq5/2$.

\section{Scope and attribution}
The source's introductory alternating-group Conjecture~1.1 is attributed
to its earlier authors. The already accessible paper \emph{Highest
rank of a polytope for $A_n$} (catalogue L044) states that it resolves
that conjecture. It is not a new solution or an open-question claim here.
The introductory symmetric-group and intransitive-product examples
are supplied by the declared common extremal theorem and the elementary
direct-product construction. Mar\'oti's theorem is retained in its
original precise form. The finer Table~3 maximum-rank reports remain attributed ancillary reports,
except the explicitly recovered entries in the clause map; none is a premise
of the principal proof.
\end{document}
